\chapter{Abstract}
In the Australian health system, Government reimbursement decisions are often made on the basis of opportunity costs. The methods used to inform these decisions are health economic evaluation and health technology assessment. These consider the costs and consequences of alternative courses of action and have not traditionally considered environmental impacts. 

There is interest in Australia and elsewhere in understanding methods to include environmental impacts in health technology assessment in the pursuit of minimising the environmental and climate change impact of health technologies. 

The aim of this thesis was to evaluate methods to include environmental impacts in health economic evaluations and health technology assessments in the Australian reimbursement decision-making context. 

Chapter 1 introduces the research and frames the aims and objectives of the thesis.

Chapter 2 outlines the methods and structure of the thesis.

Chapter 3 shows environmental impacts can be estimated alongside costs in a cost analysis and demonstrates this using the same economic methods that identify, measure, and value resources.

Chapter 4 identifies at least four approaches and ten methods to include environmental impacts in health economic evaluations, many of which have not been demonstrated with empirical evidence. 

Chapter 5 estimates the health system costs and carbon emissions of diagnostic imaging by patients treated for stage III melanoma on different PET/CT surveillance schedules. 

Chapter 6 estimates the carbon footprint of full skin examination, total body photography, and the excision of a lesion suspicious for melanoma using process-based life cycle assessment methods.  

Chapter 7 demonstrated parallel and integrated methods to include carbon emissions in a cost-effectiveness analysis of 12-monthly PET/CT imaging for patients treated for stage III melanoma.

Chapter 8 describes items to include in a health economic analysis plan (HEAP) to describe the planned economic evaluation of a randomised controlled trial. 

Chapter 9 discusses methods to include environmental impacts in health technology assessment in the Australian policy context and proposes a simple decision matrix for the cost-minimisation analysis pathway.

Chapter 10 discusses the findings of the thesis.

The thesis found that economic evaluation frameworks are a useful way to estimate and present the environmental impact of health technologies. Further work is needed to understand how these can most usefully be implemented to inform HTA decisions in Australia and globally. 